\section{Set codes and models indexed by orders}
\label{found:section-reduction-orders}

We first separate the uniformity over AECs from the transfer argument for
one AEC. All directed families below are nonempty and set-indexed. We
begin by deriving the directed forms of the chain axioms in
Definition~\ref{main:aec}.

\begin{lemma}[Directed unions and coherent direct limits]
\label{found:directed-colimits}
If a nonempty set of models is directed under strong inclusion, its
union belongs to $\K$ and every member is strong in the union. If
every member is strong in a common model $N$, its union is strong in
$N$. More generally, a coherent diagram of strong maps indexed by a
nonempty directed partial order has a direct limit in $\K$, with
strong canonical maps. Every compatible family of strong maps from
the diagram into a model induces a strong map from the direct limit.
\end{lemma}

\begin{proof}
The union is a well-defined $\tau$-structure: finitely many
arguments are contained in one family member, and any two
choices of such a member agree after passing to a common
strong upper member. Here finitarity is used for functions
and relations alike.
We prove the union and common-upper-bound assertions together by
induction on the cardinality of the family. A finite directed
family has a member containing every other member, so its union is
that member. A countable directed family has a cofinal increasing
sequence: enumerate the family and successively choose a member
above the previous choice and the next enumerated member. The
chain axioms apply to this sequence, including their
common-upper-bound clause.

Suppose the family is indexed by a set $I$ of uncountable
cardinality $\kappa$, and the two assertions hold for smaller
families. Choose an operation assigning to each pair of indices
an index of a strong upper bound. Express $I$ as the union of an
increasing continuous sequence $(I_\alpha)_{\alpha<\cf(\kappa)}$
of nonempty subsets of size less than $\kappa$, each closed under
that operation. To obtain this sequence, enumerate $I$ by
$\kappa$, use a cofinal sequence of ordinals below $\kappa$
to specify the required initial segments, close each stage
under the binary operation by countably many iterations, and
take unions at limit stages. A union of fewer than
$\cf(\kappa)$ sets of size less than $\kappa$ still has size
less than $\kappa$; countable closure preserves the size of
each infinite stage. This construction also works when
$\kappa$ is a successor cardinal, using $\kappa$ many stages.

Let $U_\alpha$ be the union of the family on $I_\alpha$. The
induction hypothesis gives $U_\alpha\in\K$. For $\alpha<\beta$,
every member indexed by $I_\alpha$ is strong in $U_\beta$.
The inductive common-upper-bound assertion therefore gives
$U_\alpha\leq_\K U_\beta$. The chain axioms now give the
required union over $I$. If the original family is strong in
$N$, induction first gives $U_\alpha\leq_\K N$ for every
$\alpha$, and the chain common-upper-bound clause gives the
same conclusion for their union. This completes the induction.

For a diagram $(M_i,f_{ij})$ indexed by a directed partial order,
form the quotient of the disjoint union of its underlying sets
by
\[
(i,x)\sim(j,y)
\quad\Longleftrightarrow\quad
(\exists k\geq i,j)\ f_{ik}(x)=f_{jk}(y).
\]
Directedness and coherence make this an equivalence relation.
The canonical map from each $M_i$ is injective. Interpret any
finitary function or relation in a common later model containing
representatives of its arguments; coherence and the fact that
the maps are structure embeddings make this interpretation
well-defined. The images of the $M_i$ are models forming a
directed family of strong inclusions. The first assertion puts
their union, which is the quotient just constructed, in $\K$
and makes all canonical maps strong.

Finally, compatible strong maps $g_i:M_i\longrightarrow N$
induce a map on that quotient. It is injective: if two images
agree, compare their representatives in a common later $M_k$
and use injectivity of $g_k$. The same argument gives a
structure embedding. Its image is the directed union of the
strong submodels $g_i[M_i]$ of $N$, so the
common-upper-bound assertion makes the induced map strong.
The quotient construction also gives uniqueness, and hence
the stated direct-limit property.
\end{proof}

\begin{lemma}[Recovery from small models and pairs]
\label{found:small-code}
Fix an infinite cardinal $\lambda$ and a finitary language $\tau$ of
cardinality at most $\lambda$. An AEC in $\tau$ with
$\LS(\K)\leq\lambda$ is determined by its models of cardinality at most
$\lambda$ and its strong inclusions between such models, both taken up to
isomorphism. These data have a set code, and the recovered class and
strong-submodel relation are uniformly first-order definable in set
theory from that code.
\end{lemma}

\begin{proof}
Represent small structures on ordinals of cardinality at most $\lambda$.
Represent a small pair $M\subseteq N$ by such a representative for $N$
together with the distinguished subset that is the domain of $M$.
The possible structures and pairs form sets. A code specifies which
structure types and which pair types are allowed. Testing whether a
given small structure or pair has an allowed type is a set-theoretic
condition: quantify over the relevant isomorphisms.

For an actual $M\in\K$, let
\[
\mathcal D(M)=\{A\leq_\K M:\|A\|\leq\lambda\}.
\]
This is a set. It is nonempty by the LS axiom applied to the empty set,
and its union is $M$ by the LS axiom applied to singletons. If
$A,B\in\mathcal D(M)$, the LS axiom inside $M$ supplies a small strong
submodel containing $|A|\cup|B|$. Coherence makes both $A$ and $B$ strong
in that submodel. Thus $\mathcal D(M)$ is directed.

Conversely, suppose a structure $M$ is the union of a nonempty directed
family of small structures whose membership and pair inclusions are
allowed by the code of $\K$. These are actual members and strong
inclusions of $\K$, so the directed-union axiom gives $M\in\K$.
Consequently the following condition recovers membership, including for
structures not already assumed to belong to $\K$:
\begin{quote}
$M$ is a union of a nonempty directed family of allowed small models,
with all its directed inclusions allowed by the small-pair code.
\end{quote}

To recover the strong relation, require a directed presentation
$N=\bigcup\mathcal D$ of this kind and a nonempty directed subfamily
$\mathcal E\subseteq\mathcal D$ with $M=\bigcup\mathcal E$.
If $M\leq_\K N$, use $\mathcal D=\mathcal D(N)$ and
$\mathcal E=\mathcal D(M)$; transitivity makes the second a subfamily
of the first. In the reverse direction, every member of $\mathcal E$
is strong in $N$ by the union clause for $\mathcal D$. The
common-upper-bound clause for $\mathcal E$ then gives $M\leq_\K N$.

Each presentation can be quantified over as a set of substructures of
the structure being decoded. All the preceding conditions are thus
first-order set-theoretic formulas with the small code as parameter.
They recover both the class and the order, as asserted.
\end{proof}

The equivalence below is the usual Hanf-number reduction; compare
\cite[Remark~1.3]{VaseyUniversalI}. We spell out the set coding to make
the uniform quantifier explicit.

\begin{proposition}[The uniform reduction]
\label{found:uniform-reduction}
The uniform assertion in Theorem~\ref{main:uniform} is equivalent to
the following qualitative assertion:
\begin{quote}
Every AEC categorical in unboundedly many cardinals is categorical in
all cardinals above some bound depending on that AEC.
\end{quote}
\end{proposition}

\begin{proof}
At a fixed LS bound $\lambda$, the languages can be renamed into a
fixed set of symbols with specified finite arities. By
Lemma~\ref{found:small-code}, the possible AECs in those languages are
determined by members of a set of small codes. A code is valid exactly
when its decoded class and relation satisfy the AEC axioms with LS
bound $\lambda$. This is a first-order condition on the code, so the
valid codes form a set.

Assume the qualitative assertion. For each valid code $c$, its
categoricity spectrum is either bounded or unbounded. In the bounded
case choose the least infinite cardinal $b_c$ strictly above every
categoricity cardinal. In the unbounded case choose the least infinite
cardinal $b_c$ such that the decoded AEC is categorical in every
cardinal at least $b_c$. These choices are definable, and the second
choice exists by the qualitative assertion. Replacement bounds the
set of all $b_c$. Choose a cardinal $\mu(\lambda)\geq\lambda$ strictly
above that set of bounds. If an AEC with LS number at most $\lambda$
is categorical at some $\kappa\geq\mu(\lambda)$, its code cannot be
in the bounded-spectrum case. It is therefore categorical at every
$\kappa'\geq b_c$, in particular at every
$\kappa'\geq\mu(\lambda)$. This is exactly the same threshold in
the premise and conclusion.

Conversely, for an AEC categorical unboundedly often, apply the uniform
assertion with $\lambda=\LS(\K)$. Some categoricity cardinal crosses
the resulting threshold, and the asserted tail follows.
\end{proof}

\begin{remark}
Lemma~\ref{found:small-code} also supplies the usual ZFC interpretation
of quantification over AECs: one quantifies over their set codes and
uses the decoded class predicates. No choice of representatives from
a proper class and no class choice axiom is needed. Below we fix one
AEC categorical unboundedly often and write $\ell=\LS(\K)$. It has
arbitrarily large models because categoricity includes existence.
\end{remark}

\subsection{A finite-arity partition argument}

The following end-homogeneous argument belongs to classical partition
calculus; compare Erd\H{o}s and Rado
\cite[Theorem~39 and its corollary]{ErdosRado1956}.  We prove the
finite-arity form and bounds used by the order-template construction.

\begin{lemma}
\label{found:partition}
For infinite cardinals $\xi,\nu$ and a positive integer $k$, there is
an infinite cardinal $R_k(\xi,\nu)$ such that every coloring of the
$k$-element subsets of an ordinal of cardinality at least
$R_k(\xi,\nu)$ into $\xi$ colors has a homogeneous subset of
cardinality $\nu$. Finitely many specified arities can be made
homogeneous simultaneously by enlarging the initial cardinal.
\end{lemma}

\begin{proof}
For $k=1$, a set of cardinality greater than $\xi\cdot\nu$ cannot be
partitioned into $\xi$ sets all of cardinality less than $\nu$.
Suppose the assertion is known for $k$. Choose an infinite initial
ordinal $\alpha$ large enough for that $k$-ary assertion with target
$\nu$, and put $b=\xi^{|\alpha|}$. Consider a coloring $c$ of the
$(k+1)$-element subsets of an ordinal $\Omega$ with $|\Omega|>b$.

Construct a tree of residual subsets of $\Omega$ through height
$\alpha$. Its root is $\Omega$. At a nonempty node, name and remove
its least point $x$. The previously named points on its branch,
together with $x$, form an increasing sequence. Partition the
remaining residual set according to all values
\[
u\longmapsto c(s\cup\{u\}),
\qquad
s\text{ a $k$-subset of that sequence}.
\]
The nonempty parts are the children. At a limit height, a branch has
as its residual set the intersection of the earlier residual sets;
retain the node precisely when this intersection is nonempty.

At each node there are at most $|\alpha|$ coordinates to record and
hence at most $b$ children. The total number of possible nodes
through height $\alpha$ is at most
$b^{|\alpha|}=b$. This includes the limit nodes, which are specified
by their branches. At most $b$ points are removed as node labels.
Choose $u\in\Omega$ which is never removed. At every height it
belongs to one residual set; at limit heights it belongs to the
intersection on that branch. Consequently it determines a branch
with $\alpha$ named points and nonempty remainder.

For every $k$-subset $s$ of these named points, all later named points
give the same value of $c(s\cup\{x\})$: that value was fixed by the
partition immediately after the last point of $s$ was named.
Therefore $c$ on increasing $(k+1)$-tuples from the branch depends
only on their first $k$ entries. Apply the $k$-ary induction
hypothesis to this induced coloring on the branch. Its homogeneous
$\nu$-subset is homogeneous for $c$ as well. For example, the
recursion just proved allows
$R_{k+1}(\xi,\nu)=(\xi^{R_k(\xi,\nu)})^+$ after increasing the
initial $R_1$ as needed.

For a finite list of arities, work backward to choose a bound large
enough for the last desired thinning after all the earlier ones.
Successive thinning preserves homogeneity already obtained. This
proves the final assertion.
\end{proof}

\subsection{A fixed presentation on every order}

The construction is an Ehrenfeucht--Mostowski presentation
\cite{EhrenfeuchtMostowski1956}; for its AEC form compare
\cite[Definition~IV.0.8 and Claims~IV.0.9, IV.0.12]{Shelah2009}.
We include the finite-support construction in the label notation used below.

\begin{theorem}[Order presentation]
\label{found:order-template}
Suppose $\K$ has arbitrarily large models and $\LS(\K)=\ell$.
There are at most $\ell$ labels, each of positive finite arity, and a
functor $E$ from nonempty linear orders and order embeddings to
$\K$ and strong maps, with the following properties.
\begin{enumerate}
\item Every element of $E(I)$ is a label evaluated on an increasing
finite list from $I$. Equality and all structure relations and
function values are determined by the labels and the order/equality
pattern of the finitely many indices involved.
\item $E(I)$ contains distinct distinguished elements $(a_i)_{i\in I}$.
An order embedding $f:I\longrightarrow J$ sends $a_i$ to $a_{f(i)}$,
and its action on any labeled element is position substitution.
\item For every finite nonempty $s\subseteq I$, the elements labeled
on $s$ form a strong submodel $E_s(I)$ of cardinality at most $\ell$.
For $s\subseteq t$ these submodels are strongly included, and
$E(I)=\bigcup_s E_s(I)$.
\item $|I|\leq\|E(I)\|\leq |I|+\ell$. In particular, for every
infinite $\kappa\geq\ell$, $\K$ has a model of cardinality exactly
$\kappa$.
\end{enumerate}
Here a subpresentation is identified with its image under the
canonical strong map when an inclusion of underlying universes is
convenient.
\end{theorem}

\begin{proof}
Take an arbitrarily large $M\in\K$ and distinct elements
$(a_i)_{i\in A}$ indexed by a large ordinal $A$. Recursively on the
positive finite size of $s\subseteq A$, choose $M_s\leq_\K M$ of
cardinality at most $\ell$ containing the displayed elements on $s$
and every $M_t$ for nonempty $t\subsetneq s$. There are only finitely
many such $t$, so the LS axiom applies with bound $\ell$. Coherence
gives $M_t\leq_\K M_s$.

Choose a surjective enumeration
$(m_{s,\eta})_{\eta<\ell}$ of each $M_s$, allowing repetitions. If
$s=\{i_0<\cdots<i_{n-1}\}$, require $m_{s,j}=a_{i_j}$ for $j<n$.
For an $n$-element set $s$, its color is the complete atomic diagram,
including negations and equalities, of all the enumerations
\[
(m_{t,\eta})_{\eta<\ell},
\qquad \varnothing\ne t\subseteq s,
\]
with the subsets of $s$ named by their relative positions. This
diagram includes formulas involving entries from different
enumerations. There are at most $\ell$ names and atomic formulas
for each fixed $n$, so its possible colors belong to a fixed set of
cardinality at most $2^\ell$.

We choose colors $d_n$, for $n\geq1$, so that for every finite $n$
there are samples of arbitrarily large cardinality homogeneous in
colors $d_1,\ldots,d_n$. Suppose $d_1,\ldots,d_{n-1}$ have been
chosen with this property. For any prescribed target cardinal,
choose such a sample large enough for Lemma~\ref{found:partition}
at arity $n$ and thin it to an $n$-homogeneous set of the target
cardinality. The earlier colors are retained. If every possible
$n$-color supported only boundedly large such thinnings, the
supremum of the bounds over the set of colors would bound all
thinnings, a contradiction. Thus one color $d_n$ supports
arbitrarily large thinnings. It restricts on every smaller subset
to the previously chosen colors. This proves the recursion and
its coherence. It asserts finite-stage realizability in
arbitrarily large samples; it does not require an infinite
simultaneous homogeneous set inside one sample.

For a nonempty linear order $I$, introduce symbols $(s,\eta)$ for
each finite nonempty $s\subseteq I$ and $\eta<\ell$, with $s$
listed increasingly. Equate $(s,\eta)$ and $(t,\zeta)$ when the
color $d_{|s\cup t|}$ says that the corresponding entries are
equal. Coherence makes this independent of passage to a larger
finite set of indices. Reflexivity, symmetry, and transitivity
hold because any finite instance occurs in a realized color.
Form the quotient by this equivalence relation.

Interpret a relation on finitely many quotient elements in a
color on the union of their supports. For a function symbol,
evaluate it in the enumerated model on that union and use a name
for its value. All its arguments belong to that enumerated model,
and the model is a substructure, so the value has such a name.
Changing representatives or enlarging the finite support gives
the same quotient element by the complete atomic diagram and
coherence. A constant can be evaluated using any one nonempty
finite support. Two such evaluations agree in a color on their
union. Thus these rules define a $\tau$-structure $E(I)$.

Let $E_s(I)$ consist of the elements named by $(s,\eta)$,
$\eta<\ell$. The color $d_{|s|}$ makes this a full enumerated copy
of one of the actual models $M_s$ in a sample. Names with support
$t\subseteq s$ are included in this copy: their equalities with
names in the enumeration on $s$ were recorded. If $s\subseteq u$,
the pair $E_s(I)\subseteq E_u(I)$ is isomorphic to an actual pair
$M_s\leq_\K M_u$ from a realization of $d_{|u|}$. Isomorphism
invariance therefore makes the inclusion strong. The family of
these small models is directed and covers the constructed
structure. Its directed union belongs to $\K$, with all
$E_s(I)\leq_\K E(I)$.

For an order embedding $f:I\longrightarrow J$, substitution
$(s,\eta)\mapsto(f[s],\eta)$ respects every defining diagram and
every inequality of names. It therefore induces an embedding
$E(f):E(I)\longrightarrow E(J)$. On each finite-support model it
is an isomorphism onto $E_{f[s]}(J)\leq_\K E(J)$. The
common-upper-bound clause applied to their directed union shows
that the whole image is strong in $E(J)$. Substitution also
proves $E(\id)=\id$ and $E(g\circ f)=E(g)\circ E(f)$.

Define $a_i$ by the first designated name on the singleton $\{i\}$.
The recorded equalities identify this with the designated name
for $i$ on any larger finite support. Distinct indices give
distinct elements because every relevant sample had that
property. The labels are indexed by $(n,\eta)$ with $n\geq1$
finite and $\eta<\ell$, so there are at most $\ell$ of them.
Counting finite supports proves the upper bound on size, and the
distinct displayed elements prove the lower bound.
\end{proof}
